\addcontentsline{toc}{chapter}{Zusammenfassung}
Eine schnelle Erdbeobachtung setzt zwei Ereignisse in der richtigen Reihenfolge voraus: Ein Satellit muss eine Aufgabe empfangen und danach vor Ablauf der Frist einen geeigneten Zugang zum Ziel erhalten. Diese Arbeit untersucht, ob ausgewaehlte Satelliten mit einer Zentral-Knoten-Rolle diese Kette in einem verteilten Satellitensystem verbessern koennen. Ein Zentral-Knoten ist hier ein bevorzugter Relay-Knoten mit erweiterter Link-Berechtigung und verbessertem Kommunikationsbudget. Er ist weder ein zentraler Missionsplaner noch ein Satellit mit einem anderen Sensor.

Aktive-Feuer-Produkte von Sentinel-3 SLSTR werden in nachvollziehbare Aufgaben fuer eine nachfolgende Beobachtung umgewandelt. Das Experiment variiert Satellitenzahl, Bahnebenen, Hoehe, Inklination und Zentral-Knoten-Anteil. Das Routing wird ueber endliche, zeitabhaengige Kommunikationsfenster simuliert. Eine Beobachtung zaehlt nur dann als erfuellt, wenn der beobachtende Satellit zuvor die vollstaendige Aufgabe empfangen hat. Die Auswertung trennt Missionsleistung, Kommunikationsaufwand, Fehlerrobustheit und Entscheidungsstabilitaet.

Die vergleichbare Analyse der 60- und 120-Satelliten-Faelle zeigt, dass die dichtere indische Aufgabenmenge auf identischen Architekturen eine um 4,66 Prozentpunkte hoehere Beobachtungserfuellung erreicht. Gleichzeitig sind die vollstaendige fristgerechte Verteilung und die maximale Verteilungszeit schlechter als in Kalifornien. Ein hoeherer Zentral-Knoten-Anteil verbessert einzelne Kennzahlen, jedoch nicht monoton. Grosse Paketmengen sind ein wichtiger Engpass: Eine Vervierfachung aller Pakete senkt die strikte Verteilung in Kalifornien um 14,94 und in Indien um 37,02 Prozentpunkte. Die Ergebnisse begruenden deshalb einen Trade-Space und kein universelles Optimum. Eine getrennte indische Sensitivitaetsanalyse mit 180 Satelliten wird nur bedingt interpretiert, da die Provenienzpruefung eine konsistente, aber unbeabsichtigte Aufgabenmenge von 741 Aufgaben ergab.

Der wesentliche Beitrag ist eine reproduzierbare Verbindung von Wildfire-Daten zu kommunikationsbewussten Architekturentscheidungen. Geometrischer Zugang allein reicht nicht aus. Routing, Arbeitslast, Link-Kapazitaet und Fristen bestimmen gemeinsam, ob aus einer moeglichen Beobachtung ein rechtzeitig abgeschlossenes Missionsergebnis wird.
